\documentclass[11pt,a4paper]{article}

% --- Codificación e idioma ---
\usepackage[utf8]{inputenc}
\usepackage[T1]{fontenc}
\usepackage[spanish,es-tabla]{babel}
\usepackage{lmodern}

% --- Geometría y formato general ---
\usepackage[margin=2.5cm]{geometry}
\usepackage{microtype}
\usepackage{parskip}
\usepackage{enumitem}
\setlist{topsep=2pt,itemsep=2pt,parsep=0pt}

% --- Tablas y diagramas ---
\usepackage{array}
\usepackage{booktabs}
\usepackage{amsmath,amssymb}
\usepackage{tikz}
\usetikzlibrary{positioning,arrows.meta,trees,calc}

% --- Cabeceras y numeración ---
\usepackage{fancyhdr}
\pagestyle{fancy}
\fancyhf{}
\fancyhead[L]{\small Bases de Datos II}
\fancyhead[R]{\small Resumen Teórico --- Proc. y Opt. de Consultas}
\fancyfoot[C]{\thepage}
\renewcommand{\headrulewidth}{0.4pt}

% --- Color y código ---
\usepackage{xcolor}
\definecolor{codebg}{RGB}{248,248,248}
\definecolor{codeframe}{RGB}{200,200,200}
\definecolor{codekw}{RGB}{0,0,180}
\definecolor{codestr}{RGB}{163,21,21}
\definecolor{codecmt}{RGB}{0,128,0}
\definecolor{notebg}{RGB}{255,250,230}
\definecolor{noteframe}{RGB}{220,180,80}

\usepackage{listings}
\lstdefinelanguage{SQLext}[]{SQL}{
    morekeywords={NUMERIC,VARCHAR,DECIMAL,DATE,TIMESTAMP,
        BIGINT,VARCHAR2,NUMBER,SCHEMA,DATABASE,
        WITH,FROM,TO,ON,AS,AND,OR,IS,NOT,NULL,
        INSERT,UPDATE,DELETE,SELECT,EXPLAIN,ANALYZE,BUFFERS,
        NATURAL,JOIN,CROSS,INNER,OUTER,LEFT,RIGHT,FULL,GROUP,BY,ORDER,
        SET,RESET,CREATE,INDEX,IF,EXISTS,VIEW,MATERIALIZED,REFERENCES},
    sensitive=false,
    morecomment=[l]{--},
    morestring=[b]'
}

\lstdefinestyle{sqlstyle}{
    language=SQLext,
    basicstyle=\ttfamily\footnotesize,
    keywordstyle=\color{codekw}\bfseries,
    stringstyle=\color{codestr},
    commentstyle=\color{codecmt}\itshape,
    backgroundcolor=\color{codebg},
    frame=single,
    rulecolor=\color{codeframe},
    breaklines=true,
    breakatwhitespace=true,
    columns=fullflexible,
    keepspaces=true,
    showstringspaces=false,
    captionpos=b,
    xleftmargin=0.4em,
    xrightmargin=0.4em,
    aboveskip=8pt,
    belowskip=8pt,
    tabsize=2
}
\lstset{style=sqlstyle}

% --- Cajas de nota ---
\usepackage[most]{tcolorbox}
\newtcolorbox{nota}{colback=notebg,colframe=noteframe,arc=2pt,
    left=6pt,right=6pt,top=4pt,bottom=4pt,boxrule=0.6pt}

% --- Hipervínculos ---
\usepackage[hidelinks,bookmarks=true]{hyperref}

% --- Comandos cortos para álgebra relacional ---
\newcommand{\sel}[1]{\sigma_{#1}}
\newcommand{\proy}[1]{\Pi_{#1}}
\newcommand{\jn}{\bowtie}

\title{Resumen Teórico --- Procesamiento y Optimización de Consultas\\[4pt]
       \large Bases de Datos II --- Práctico N.\textsuperscript{o} 5}
\author{Tomás Rodeghiero}
\date{2026}

\begin{document}

\begin{titlepage}
    \centering
    \vspace*{1.5cm}

    {\large \textbf{Universidad Nacional de Río Cuarto}}\\[6pt]
    {\normalsize Facultad de Ciencias Exactas, Físico-Químicas y Naturales}\\[4pt]
    {\normalsize Departamento de Computación}\\[4pt]
    {\normalsize Licenciatura en Ciencias de la Computación}\\[3cm]

    {\Large \textbf{Bases de Datos II}}\\[1.2cm]

    {\Large \textbf{Resumen Teórico}}\\[6pt]
    {\large Procesamiento y Optimización de Consultas\\[2pt]
    Práctico N.\textsuperscript{o} 5}\\[3cm]

    \begin{tabular}{rl}
        \textbf{Alumno:} & Tomás Rodeghiero \\[4pt]
        \textbf{Año:}    & 2026 \\
    \end{tabular}

    \vfill
    {\small Síntesis de los Teóricos 6 (Procesamiento de Consultas) y 7
    (Optimización de Consultas).}
\end{titlepage}

\tableofcontents
\newpage

% =====================================================
\section{Visión general}
% =====================================================

Cuando un motor de base de datos recibe una consulta SQL no la ejecuta
literalmente: la procesa en \textbf{tres etapas}, cada una con responsabilidades
bien delimitadas:

\begin{enumerate}
    \item \textbf{Parsing y traducción.} Se verifica la sintaxis y que las
    relaciones y columnas referenciadas existan. Luego el SQL se traduce a
    una expresión del \emph{álgebra relacional}. Si la consulta usa
    \emph{vistas}, la traducción reemplaza la vista por su definición.
    \item \textbf{Optimización.} A partir de la expresión inicial se generan
    \emph{expresiones equivalentes} (mismo resultado, distinto orden o forma)
    y, por cada una, varios \emph{planes de ejecución} (qué algoritmo concreto
    se usa para cada operador). El optimizador estima costos y elige el plan
    más barato.
    \item \textbf{Evaluación.} El motor toma el plan elegido y lo ejecuta
    contra los datos en disco, devolviendo el resultado.
\end{enumerate}

\begin{center}
\begin{tikzpicture}[node distance=8mm, every node/.style={font=\footnotesize}]
    \node[draw,rounded corners] (q) {\strut Consulta SQL};
    \node[draw,rounded corners,right=10mm of q] (p) {\strut Parser y traductor};
    \node[draw,rounded corners,right=10mm of p] (e) {\strut Expresión $\sigma,\Pi,\bowtie$};
    \node[draw,rounded corners,below=8mm of e] (o) {\strut Optimizador};
    \node[draw,rounded corners,below=8mm of p] (m) {\strut Motor de ejecución};
    \node[draw,rounded corners,left=10mm of m] (r) {\strut Resultado};
    \node[draw,cylinder,shape border rotate=90,aspect=0.25,right=10mm of o] (s) {Estadísticas};

    \draw[-{Latex}] (q) -- (p);
    \draw[-{Latex}] (p) -- (e);
    \draw[-{Latex}] (e) -- (o);
    \draw[-{Latex}] (s) -- (o);
    \draw[-{Latex}] (o) -- node[right,font=\tiny]{plan de ejec.} (m);
    \draw[-{Latex}] (m) -- (r);
\end{tikzpicture}
\end{center}

El \emph{punto importante} es que la diferencia de costo entre dos planes
equivalentes puede ser enorme (varios órdenes de magnitud). Por eso la
optimización no es un lujo: es lo que hace que SQL sea usable a escala.

% =====================================================
\section{Costos: cómo se mide la eficiencia}
% =====================================================

El costo de una consulta es básicamente el tiempo de respuesta. En la práctica
intervienen muchos factores (CPU, RAM, red, disco), pero el Teórico~6
simplifica: \textbf{el costo dominante es el acceso a disco}, y se calcula
contando dos cosas:

\begin{itemize}
    \item número de \emph{bloques transferidos} (lectura/escritura), con
    tiempo unitario $t_T$;
    \item número de \emph{búsquedas} (\emph{seeks}, posicionamientos del
    cabezal o cambio de partición), con tiempo unitario $t_b$.
\end{itemize}

Si una operación transfiere $b$ bloques y hace $s$ búsquedas, su costo
estimado es

\[
\text{Costo} \;=\; b \cdot t_T \;+\; s \cdot t_b .
\]

\paragraph{Valores de referencia (Teórico 6).}
$t_b = 4$~ms, $t_T = 0{,}1$~ms, bloque de 4~KB, velocidad de transferencia
40~MB/s. Son valores ``académicos''; en SSD modernos $t_b$ es mucho menor
($t_b \approx 0$ en NVMe), pero el modelo sigue siendo correcto: \emph{las
búsquedas son caras frente a las transferencias contiguas}.

\paragraph{Variables del catálogo que aparecen en los costos.}
\begin{itemize}
    \item $n_r$: cantidad de tuplas de la relación $r$.
    \item $b_r$: cantidad de bloques que contienen tuplas de $r$.
    \item $t_r$: tamaño promedio de una tupla.
    \item $f_r$: \emph{factor de bloque}, tuplas por bloque ($b_r = \lceil n_r/f_r\rceil$).
    \item $V(A,r)$: cantidad de valores distintos del atributo $A$ en $r$.
    \item $h_i$: altura del índice (B+tree).
\end{itemize}

Estas variables son las mismas que cualquier motor real mantiene en su
catálogo (en PostgreSQL: \texttt{pg\_class.reltuples}, \texttt{relpages},
\texttt{pg\_stats.n\_distinct}, etc.).

% =====================================================
\section{Algoritmos para la operación de selección}
% =====================================================

La selección $\sel{\theta}(r)$ es el operador más estudiado porque interviene
en casi todas las consultas. El Teórico~6 lista once algoritmos identificados
como A1\,\ldots\,A11; el optimizador elige uno según haya o no haya índices
y según la forma de $\theta$ (igualdad, rango, conjunción, disyunción).

\subsection{Selecciones simples (un solo predicado)}

\subsubsection*{A1: búsqueda lineal}

Lee bloque por bloque, descarta los que no cumplen el predicado. Es el
``algoritmo universal'': siempre se puede aplicar.

\[
\text{Costo}_{A1} = b_r \cdot t_T + t_b.
\]

Si la condición es de \emph{igualdad por clave} (a lo sumo una tupla
cumple), la búsqueda se corta en cuanto aparece: en promedio
$(b_r/2)\,t_T + t_b$.

\paragraph{Por qué importa.} A1 es el plan de referencia. Cualquier otro
plan vale la pena sólo si es más barato que A1. Por eso, cuando la
selectividad es muy alta, los motores eligen Seq Scan aunque haya índice
disponible: leer todo en orden es más barato que saltar de bloque en
bloque.

\subsubsection*{A2: búsqueda binaria}

Aplicable si el archivo está \emph{ordenado} por el atributo y la condición
es de igualdad: localiza la primera tupla con costo
$\lceil \log_2 b_r \rceil \cdot (t_T + t_b)$ y luego transfiere los bloques
contiguos.

\subsubsection*{A3, A4, A5: igualdad con índice}

\begin{itemize}
    \item \textbf{A3} --- igualdad sobre \emph{índice de clave primaria o
    único}. A lo sumo una tupla. Costo $(h_i + 1)(t_T + t_b)$.
    \item \textbf{A4} --- igualdad sobre \emph{índice primario no único}
    (atributo no clave, pero el archivo está ordenado por él). Las tuplas
    que cumplen están en bloques contiguos:
    $h_i\,(t_T + t_b) + t_b + b \cdot t_T$, con $b$ = bloques que
    contienen tuplas que cumplen.
    \item \textbf{A5} --- igualdad sobre \emph{índice secundario}. Si el
    atributo no es único pueden recuperarse varias tuplas, una I/O por
    tupla en el peor caso: $(h_i + n)\,(t_T + t_b)$.
\end{itemize}

\paragraph{Intuición.} A3 es el caso ideal (devuelve una sola tupla y baja
el árbol del índice una sola vez). A5 paga por cada tupla devuelta porque
puede estar en un bloque distinto: si la selectividad es muy baja, A5 puede
volverse peor que A1.

\subsubsection*{A6 y A7: rango con índice}

\begin{itemize}
    \item \textbf{A6} --- rango sobre \emph{índice primario}. Para $A\ge v$
    se busca la primera tupla con el índice y se sigue secuencialmente.
    Para $A\le v$ se arranca desde el inicio del archivo y se para en la
    primera tupla que viole la condición; el índice no agrega valor en este
    último caso.
    \item \textbf{A7} --- rango sobre \emph{índice secundario}. Se recorren
    las hojas del índice para conseguir los punteros y, por cada puntero,
    se hace una I/O.
\end{itemize}

\begin{nota}
\textbf{Cuándo conviene un índice secundario sobre uno primario.} El índice
primario garantiza que las tuplas devueltas vienen en bloques contiguos. El
secundario, no: cada tupla puede estar en un bloque distinto y obliga a
\emph{tantas} I/O como tuplas, lo que rápidamente se vuelve más caro que
leer el archivo entero (A1). Por eso, para rangos amplios, los motores
suelen preferir Seq Scan al \emph{Index Scan} con índice secundario.
\end{nota}

\subsection{Selecciones complejas (conjunción y disyunción)}

\subsubsection*{A8: conjunción usando un solo índice}

Para $\sel{\theta_1 \wedge \theta_2 \wedge \cdots \wedge \theta_n}(r)$:

\begin{enumerate}
    \item Elegir la condición $\theta_i$ que dé el plan más barato según
    A1--A7.
    \item Recuperar las tuplas que cumplen $\theta_i$.
    \item Testear el resto de las condiciones en memoria.
\end{enumerate}

Es la estrategia ``elegir la condición más restrictiva''.

\subsubsection*{A9: índice multiclave}

Si existe un índice compuesto que cubre exactamente la combinación de
predicados, se usa directamente. Ejemplo: índice \texttt{(a,b)} para
$\sel{a=v_1\wedge b=v_2}(r)$.

\subsubsection*{A10: intersección de identificadores}

Para conjunciones donde \emph{cada} condición tiene un índice utilizable:

\begin{enumerate}
    \item Obtener el conjunto de punteros $P_i$ que cumple $\theta_i$ usando
    su índice (lógica A2--A7).
    \item Calcular $P_1 \cap P_2 \cap \cdots \cap P_n$.
    \item Recuperar sólo las tuplas referenciadas por la intersección.
\end{enumerate}

Es el algoritmo subyacente a los \emph{Bitmap Index Scan} de PostgreSQL.

\subsubsection*{A11: unión de identificadores (disyunción)}

Para $\sel{\theta_1 \vee \theta_2 \vee \cdots \vee \theta_n}(r)$ con un
índice para cada $\theta_i$:

\begin{enumerate}
    \item Obtener $P_i$ con el índice de $\theta_i$.
    \item Calcular $P_1 \cup P_2 \cup \cdots \cup P_n$ (sin duplicados).
    \item Recuperar las tuplas referenciadas.
\end{enumerate}

Si alguna condición no tiene índice, no aplica y hay que caer a A1.

\paragraph{Resumen de la familia A1--A11.}

\begin{center}
\begin{tabular}{lll}
\toprule
Algoritmo & Aplicabilidad & Idea \\
\midrule
A1  & Siempre                            & Búsqueda lineal \\
A2  & Archivo ordenado, igualdad         & Búsqueda binaria \\
A3  & Igualdad, índice único             & Una sola tupla \\
A4  & Igualdad, índice primario no único & Tuplas contiguas \\
A5  & Igualdad, índice secundario        & I/O por tupla \\
A6  & Rango, índice primario             & Recorrido secuencial \\
A7  & Rango, índice secundario           & Recorrer hojas + I/O \\
A8  & Conjunción + un índice             & Filtrar el resto en RAM \\
A9  & Conjunción + índice multiclave     & Un solo acceso \\
A10 & Conjunción + varios índices        & Intersección de RIDs \\
A11 & Disyunción + índices               & Unión de RIDs \\
\bottomrule
\end{tabular}
\end{center}

% =====================================================
\section{Ordenamiento}
% =====================================================

\paragraph{Por qué importa.} Muchas operaciones (group by, distinct, merge
join, order by) requieren entrada ordenada. Cómo se ordena depende de si
la relación entra en memoria:

\begin{itemize}
    \item Si la relación cabe en RAM: cualquier algoritmo clásico
    (\emph{quicksort}, \emph{heapsort}, \ldots).
    \item Si no cabe: \textbf{External Sort-Merge}.
    \item Si ya existe un índice sobre el atributo a ordenar: alcanza con
    recorrerlo, una I/O por tupla. Conviene sólo si la cantidad de tuplas
    es chica (de lo contrario es más barato ordenar de cero).
\end{itemize}

\paragraph{External Sort-Merge.} Con $M$ bloques de memoria disponibles:

\begin{enumerate}
    \item \textbf{Fase de creación de corridas}: leer $M$ bloques a la vez,
    ordenarlos en memoria y escribir cada bloque ordenado como una
    \emph{corrida} $S_i$ a disco. Quedan $N = \lceil b_r / M\rceil$ corridas
    ordenadas pero no entre sí.
    \item \textbf{Fase de merge $N$-way}: si $N < M$, se mezclan las $N$
    corridas con $N$ buffers de entrada y 1 de salida. Si $N \ge M$, se
    hacen varias pasadas, mezclando $M-1$ corridas en cada una y reduciendo
    el número total por ese factor.
\end{enumerate}

% =====================================================
\section{Algoritmos para la operación de join}
% =====================================================

El join $r \jn_\theta s$ es la operación con más algoritmos a elegir
porque es la más cara y la que más se beneficia de la elección correcta.

\subsection{Nested-loop join}

\begin{center}
\begin{tabular}{l}
\textbf{for each} tupla $t_r$ en $r$ \textbf{do} \\
\quad \textbf{for each} tupla $t_s$ en $s$ \textbf{do} \\
\quad\quad si $\theta(t_r,t_s)$ es verdadero, emitir $t_r \cdot t_s$
\end{tabular}
\end{center}

Es el más simple y se aplica a cualquier condición $\theta$ (no sólo
igualdad). Costo del peor caso:

\[
n_r \cdot b_s + b_r \text{ bloques transferidos} \;+\; n_r + b_r \text{ búsquedas}.
\]

\paragraph{Ejemplo del teórico.} Con $n_{cli}=10\,000$, $b_{cli}=400$,
$n_{fac}=5\,000$, $b_{fac}=100$:
\begin{itemize}
    \item \texttt{factura} como externa: $5000\cdot 400 + 100 = 2{,}000{,}100$
    bloques transferidos.
    \item \texttt{cliente} como externa: $10000\cdot 100 + 400 = 1{,}000{,}400$.
\end{itemize}

Conclusión: poner como \emph{externa} la relación que minimice el producto
$n_{ext} \cdot b_{int}$; idealmente, la más chica como interna entera en
memoria.

\paragraph{Block nested-loop.} Variante que carga \emph{bloques} de $r$
en lugar de tuplas individuales y compara contra todas las de $s$. Reduce
muchísimo el número de búsquedas. Si la relación más chica entra en memoria
el costo cae a $b_r + b_s$ con 2 búsquedas.

\subsection{Indexed nested-loop join}

Si existe un índice sobre el atributo del join en la relación interna $s$,
por cada tupla de $r$ se busca directamente en el índice. Útil para joins
selectivos o cuando $s$ es grande pero el predicado de igualdad recupera
pocas tuplas.

\subsection{Merge-join}

Aplicable a \emph{equi-joins} y joins naturales. Funciona así:

\begin{enumerate}
    \item Ordenar ambas relaciones por los atributos del join (si todavía no
    lo están).
    \item Recorrer las dos en paralelo como un \emph{merge}: avanzar la
    relación cuyo valor sea menor; cuando coinciden, emitir el cruce de
    tuplas con ese valor.
\end{enumerate}

Cada relación se recorre \emph{una vez}. Si las relaciones ya están
ordenadas (por ejemplo, por índice clusterizado o por una operación previa
que dejó el orden), Merge-Join es muy barato.

\subsection{Hash-join}

Aplicable a joins de igualdad y naturales:

\begin{enumerate}
    \item Elegir una función hash $h$ sobre los atributos del join.
    \item Particionar $r$ en buckets $r_0, r_1, \ldots, r_n$ según $h$ y
    análogamente $s$ en $s_0, \ldots, s_n$. La clave: tuplas de $r_i$ sólo
    pueden cruzarse con tuplas de $s_i$ (porque tienen el mismo hash).
    \item Por cada par $(r_i, s_i)$, hacer el join en memoria (típicamente
    \emph{nested loop} sobre la partición chica con tabla hash sobre la
    otra).
\end{enumerate}

Es el algoritmo que más usan los motores modernos cuando ninguna relación
viene ordenada. Es el caso por defecto cuando se hace
\texttt{NATURAL JOIN} sin índices ni \texttt{ORDER BY} (como pasa en el
Ejercicio~5.e).

\paragraph{Comparación práctica.}

\begin{center}
\begin{tabular}{lll}
\toprule
Algoritmo & Condiciones soportadas & Cuándo conviene \\
\midrule
Nested-loop          & Cualquiera             & Una relación muy chica o como fallback \\
Nested-loop (bloques)& Cualquiera             & Relación chica cabe en RAM \\
Indexed nested-loop  & Cualquiera + índice    & Join muy selectivo en la interna \\
Merge-join           & Igualdad / natural     & Entradas ya ordenadas \\
Hash-join            & Igualdad / natural     & Cualquier tamaño sin orden previo \\
\bottomrule
\end{tabular}
\end{center}

% =====================================================
\section{Evaluación de expresiones: materialización vs pipelining}
% =====================================================

Hasta aquí se vieron operadores individuales. Una consulta real es un árbol
de operadores. Hay dos estrategias para combinarlos:

\paragraph{Materialización.} Cada subexpresión se evalúa por completo y
su resultado se escribe en disco; el operador siguiente lee desde ese
resultado intermedio. Es simple pero paga el costo de escribir y leer la
relación intermedia entera.

\paragraph{Pipelining.} Las tuplas que va produciendo un operador se pasan
inmediatamente al operador siguiente, sin escribir el intermedio en disco.
Es como un \emph{stream} de tuplas. Es lo que hacen los motores modernos
cuando el operador siguiente es \emph{tuple-at-a-time} (selección,
proyección, nested-loop, hash-join probe).

\paragraph{Quién bloquea el pipeline.} Hay operadores que necesitan ver
\emph{toda} su entrada antes de producir el primer resultado: por ejemplo,
sort, hash-join en su fase de \emph{build}, agrupación. Estos son los que
fuerzan materialización en el árbol.

% =====================================================
\section{Optimización: generación de expresiones equivalentes}
% =====================================================

El optimizador toma la expresión inicial (resultado del traductor) y la
reescribe usando \textbf{reglas de equivalencia}: dos expresiones son
equivalentes si producen el mismo resultado (en SQL, el mismo \emph{multiset})
para cualquier instancia de la base de datos.

\subsection{Reglas más usadas (Teórico 7)}

\begin{enumerate}
    \item \textbf{R1 --- Cascada de selecciones.}
    $\sel{\theta_1\wedge\theta_2}(E) = \sel{\theta_1}(\sel{\theta_2}(E))$.
    \item \textbf{R2 --- Conmutatividad de selecciones.}
    $\sel{\theta_1}(\sel{\theta_2}(E)) = \sel{\theta_2}(\sel{\theta_1}(E))$.
    \item \textbf{R3 --- Cascada de proyecciones.}
    Sólo importa la proyección externa: $\proy{L_1}(\proy{L_2}(\cdots(E))) = \proy{L_1}(E)$.
    \item \textbf{R4 --- Selección sobre producto/theta-join.}
    $\sel{\theta}(E_1\times E_2) = E_1 \jn_\theta E_2$.
    \item \textbf{R5 --- Conmutatividad de join.}
    $E_1 \jn E_2 = E_2 \jn E_1$.
    \item \textbf{R6 --- Asociatividad de join.}
    $(E_1 \jn E_2) \jn E_3 = E_1 \jn (E_2 \jn E_3)$.
    \item \textbf{R7 --- Push-down de selección sobre join.}
    Si $\theta_0$ involucra sólo atributos de $E_1$:
    $\sel{\theta_0}(E_1\jn E_2) = \sel{\theta_0}(E_1)\jn E_2$.
    Si $\theta_1$ usa sólo $E_1$ y $\theta_2$ sólo $E_2$:
    $\sel{\theta_1\wedge\theta_2}(E_1\jn E_2) = \sel{\theta_1}(E_1)\jn \sel{\theta_2}(E_2)$.
    \item \textbf{R8 --- Push-down de proyección sobre join.}
    $\proy{L_1\cup L_2}(E_1\jn_\theta E_2)$ se distribuye, conservando en cada
    operando los atributos que pertenezcan a $L_1\cup L_2$ o aparezcan en
    $\theta$.
    \item \textbf{R9--R12 --- Operaciones de conjuntos.} Unión e intersección
    son conmutativas y asociativas; la resta no. La selección distribuye
    sobre unión, intersección y resta.
\end{enumerate}

\subsection{Intuición y ejemplo}

\paragraph{Por qué empujar selecciones.} Si una selección elimina el 99\%
de las tuplas, conviene aplicarla antes de cualquier join porque achica el
tamaño del operando que entra al join. El join trabaja sobre menos tuplas
y produce menos tuplas.

\paragraph{Ejemplo (Teórico 7).}

\[
\proy{desc}\!\big(\sel{nombre=\text{``Super Vea''}\,\wedge\,precio\_venta<10}
(proveedor \jn provee \jn articulo)\big)
\]

Aplicando R1 (cascada), R7 (push-down) y R5/R6 (reordenar) se llega a:

\[
\proy{desc}\!\big( \sel{nombre=\text{``Super Vea''}}(proveedor) \jn
\sel{precio\_venta<10}(provee) \jn articulo \big)
\]

Ahora cada selección actúa sobre la tabla base más chica y los joins
trabajan sobre operandos ya reducidos.

\subsection{Orden de joins}

Para tres relaciones, asociatividad da dos planes:
$(r_1\jn r_2)\jn r_3$ y $r_1\jn(r_2\jn r_3)$. Si $r_2\jn r_3$ resulta
mucho mayor que $r_1\jn r_2$, el primer plan es preferible (resultado
intermedio más chico). Para $n$ relaciones hay
$\dfrac{(2(n-1))!}{(n-1)!}$ órdenes posibles; con $n=10$, casi 18 mil
millones. Por eso el optimizador usa \textbf{programación dinámica}: calcula
el mejor plan para cada subconjunto de relaciones una sola vez y los
combina.

% =====================================================
\section{Optimización basada en costos vs heurística}
% =====================================================

Los optimizadores reales combinan dos enfoques:

\paragraph{Optimización basada en costos.} Enumerar (parte de) las
expresiones equivalentes posibles, estimar el costo de cada plan usando
las estadísticas del catálogo y elegir el más barato. Es preciso pero caro
(espacio de planes enorme).

\paragraph{Optimización heurística.} Aplicar reglas que \emph{típicamente}
mejoran el plan sin estimar costos para cada caso. Las heurísticas
clásicas son:

\begin{itemize}
    \item realizar las selecciones tan pronto como sea posible (reducir
    cantidad de tuplas),
    \item realizar las proyecciones tan pronto como sea posible (reducir
    ancho de tupla),
    \item realizar primero los joins/selecciones más restrictivos (los que
    den resultados intermedios más chicos).
\end{itemize}

Los motores reales hacen primero una pasada heurística (push-down,
reescritura de vistas, eliminación de subqueries) y luego una pasada
basada en costos sobre el espacio reducido.

\begin{nota}
\textbf{Atención.} Las heurísticas mejoran ``en general'', no siempre.
Empujar una selección puede ser contraproducente si la condición es muy
costosa de evaluar (por ejemplo, una función definida por usuario), o si
le quita ordenamiento útil al operador siguiente.
\end{nota}

% =====================================================
\section{Estadísticas para estimación de costos}
% =====================================================

Para que la optimización por costos funcione, el optimizador necesita
\emph{estimar} cuánto cuesta cada plan sin ejecutarlo. Esto se basa
enteramente en las estadísticas del catálogo.

\subsection{Variables que mantiene el catálogo}

\begin{itemize}
    \item $n_r$, $b_r$, $t_r$, $f_r$ por relación.
    \item $V(A,r)$: cantidad de valores distintos del atributo $A$ en $r$.
    \item $\min(A,r)$ y $\max(A,r)$.
    \item \textbf{Histogramas} sobre los atributos más consultados,
    típicamente equi-depth: dividen el rango de valores en buckets con
    igual cantidad de tuplas. Sirven para estimar selectividad de rangos
    cuando los valores no están distribuidos uniformemente.
    \item \textbf{Most Common Values (MCV)}: lista de valores más
    frecuentes con su frecuencia. Sirven para estimar igualdades cuando
    el valor pedido aparece en la lista.
\end{itemize}

\subsection{Estimación de la selección}

\paragraph{Igualdad: $\sel{A=v}(r)$.}
Si $A$ no es clave: $n_r / V(A,r)$ tuplas (asume distribución uniforme).
Si $A$ es clave: a lo sumo 1 tupla.
Si hay MCV y $v \in$ MCV: usar la frecuencia exacta.

\paragraph{Rango: $\sel{A\le v}(r)$.}
Si no hay histograma: aproximar linealmente entre $\min$ y $\max$:
\[
c \;=\; n_r \cdot \frac{v - \min(A,r)}{\max(A,r) - \min(A,r)}.
\]
Con histograma equi-depth: sumar bucket por bucket, interpolando el
bucket final.

\paragraph{Conjunción: $\sel{\theta_1\wedge\cdots\wedge\theta_n}(r)$.}
Asumiendo \emph{independencia} entre las condiciones:
\[
n_r \cdot \frac{s_1 \cdot s_2 \cdots s_n}{n_r^{\,n}},
\]
donde $s_i = |\sel{\theta_i}(r)|$. La probabilidad $s_i / n_r$ se llama
\emph{selectividad} de $\theta_i$.

\paragraph{Disyunción: $\sel{\theta_1\vee\cdots\vee\theta_n}(r)$.}
Por inclusión--exclusión, asumiendo independencia:
\[
n_r \cdot \big(1 - (1 - s_1/n_r)(1 - s_2/n_r)\cdots(1 - s_n/n_r)\big).
\]

\paragraph{Negación: $\sel{\neg\theta}(r)$.}
$n_r - |\sel{\theta}(r)|$.

\subsection{Estimación del tamaño de un join}

\paragraph{Producto cartesiano.}
$|r\times s| = n_r \cdot n_s$.

\paragraph{Si $R \cap S = \varnothing$.}
$r \jn s$ degenera en $r \times s$: $|r\jn s| = n_r \cdot n_s$.

\paragraph{Si el atributo común es clave en una de las relaciones.}
Cada tupla de la otra matchea con a lo sumo una:
$|r \jn s| \le \min(n_r, n_s)$.

\paragraph{Si el atributo común no es clave en ninguna de las dos.}
Dos estimaciones posibles, dependiendo del lado de la división:
\[
\frac{n_r \cdot n_s}{V(A,r)} \quad\text{o}\quad \frac{n_r \cdot n_s}{V(A,s)} .
\]
Se toma la \emph{menor} como mejor estimación. Si hay histogramas se aplica
la misma fórmula por celda.

\subsection{Estimación de otras operaciones}

\begin{itemize}
    \item Proyección $\proy{A}(r)$: tamaño $\le V(A,r)$ (con eliminación de
    duplicados).
    \item Unión $r\cup s$: cota superior $|r| + |s|$.
    \item Intersección $r\cap s$: cota superior $\min(|r|,|s|)$.
    \item Diferencia $r - s$: cota superior $|r|$.
\end{itemize}

% =====================================================
\section{Técnicas adicionales de optimización}
% =====================================================

\subsection{Subqueries anidados}

SQL trata conceptualmente a un subquery del \texttt{where} como una función
que se evalúa una vez por cada tupla del nivel externo (\emph{evaluación
correlacionada}). En el peor caso esto multiplica el costo por $n_r$.

Los optimizadores intentan reescribir los subqueries como joins:

\begin{lstlisting}
-- Forma anidada
SELECT descripcion FROM articulo
WHERE EXISTS (SELECT 1 FROM provee
              WHERE provee.nArt = articulo.nArt);

-- Forma equivalente con semijoin / join (mejor plan)
SELECT DISTINCT a.descripcion
FROM articulo a JOIN provee p ON p.nArt = a.nArt;
\end{lstlisting}

La forma con join expone más estructura al optimizador (puede usar hash o
merge join, índices, etc.).

\subsection{Vistas materializadas}

Una \textbf{vista materializada} es una vista cuyo contenido se calcula y
se almacena físicamente. Mejora la performance de consultas costosas a
costa de mantener actualizado el resultado.

\begin{lstlisting}
CREATE VIEW precioMenor (nArt, precio_venta) AS
SELECT nArt, min(precio_venta)
FROM provee
GROUP BY nArt;
\end{lstlisting}

\paragraph{Mantenimiento.} Tres estrategias:
\begin{itemize}
    \item Recomputar cada cierto tiempo (\emph{batch}, p.ej. cada noche).
    \item Recomputar al \texttt{commit} (\emph{refresh fast/on commit} en
    Oracle).
    \item \textbf{Mantenimiento incremental}: aplicar los cambios sobre la
    vista a medida que cambian las tablas base, usando \emph{triggers} o
    logs de cambio. Es la opción más eficiente pero la más compleja.
\end{itemize}

% =====================================================
\section{Cómo se conecta con PostgreSQL (Práctico 5)}
% =====================================================

Todo lo anterior se vuelve operativo en el Ejercicio~5 del práctico:

\begin{itemize}
    \item \textbf{Catálogo}: \texttt{pg\_class}, \texttt{pg\_attribute},
    \texttt{pg\_stats} contienen $n_r$, $b_r$, $V(A,r)$, histogramas y MCVs
    de las tablas \texttt{persona}, \texttt{recurso} y
    \texttt{seguimiento\_acceso}.
    \item \textbf{Algoritmos de selección}: el plan inicial sobre
    \texttt{seguimiento\_acceso} filtrando por \texttt{fecha\_yhora\_entrada}
    es un \emph{Seq Scan} (A1). Al crear un B+tree sobre la columna y correr
    \texttt{ANALYZE}, el plan cambia a \emph{Bitmap Index Scan + Heap Scan}
    (variantes de A5/A7 según la selectividad).
    \item \textbf{Algoritmos de join}: \texttt{NATURAL JOIN} sin orden
    previo arranca con \emph{Hash Join}. Creando un índice sobre la columna
    del join y agregando \texttt{ORDER BY} por la clave de join, el planner
    cambia a \emph{Merge Join}. Si la decisión no cambia se puede usar
    \texttt{SET enable\_hashjoin = off} para verificación académica.
    \item \textbf{Estimación}: el comando \texttt{EXPLAIN (ANALYZE, BUFFERS)}
    muestra junto al plan el costo \emph{estimado} y el costo \emph{real},
    permitiendo comparar las estimaciones del optimizador con la ejecución
    efectiva.
\end{itemize}

% =====================================================
\section{Conclusión}
% =====================================================

Procesar una consulta SQL es, en el fondo, un problema de búsqueda en un
espacio de planes equivalentes. El motor traduce el SQL a álgebra relacional,
genera planes con reglas de equivalencia, estima cada uno con las estadísticas
del catálogo y ejecuta el más barato.

Las ideas centrales para retener son:

\begin{itemize}
    \item el costo dominante es el acceso a disco (transferencias y
    búsquedas);
    \item los algoritmos de selección (A1--A11) y de join se eligen según
    los índices disponibles y la selectividad;
    \item las reglas de equivalencia permiten reescribir la expresión sin
    cambiar su resultado, y suelen mejorar el costo cuando empujan
    selecciones y proyecciones a las hojas;
    \item las estadísticas del catálogo son el insumo crítico: sin ellas,
    el optimizador estima mal y elige peor;
    \item ningún algoritmo es mejor siempre: la decisión depende de la
    forma del predicado, del tamaño relativo de las relaciones, de los
    índices y del estado de la memoria.
\end{itemize}

\end{document}
